% Chapter 3

\chapter{Methodology}

\section{Reinforcement learning formulation}

This synthetic chapter demonstrates how the template handles mathematical notation, displayed equations, aligned derivations, and multi-line optimisation objectives. The example frames a learning problem as a Markov decision process:

\begin{equation}
\mathcal{M} = \left(\mathcal{S}, \mathcal{A}, P, R, \gamma\right),
\label{eq:mdp}
\end{equation}

where $\mathcal{S}$ is the state space, $\mathcal{A}$ is the action space, $P(s' \mid s,a)$ is the transition probability, $R(s,a)$ is the expected reward, and $\gamma \in [0,1)$ is the discount factor. At each time step $t$, an agent observes a state $s_t$, selects an action $a_t \sim \pi_\theta(\cdot \mid s_t)$, receives reward $r_t$, and transitions to a new state $s_{t+1}$.

\section{Return and value functions}

The discounted return from time $t$ is defined as:

\begin{equation}
G_t = \sum_{k=0}^{\infty} \gamma^k r_{t+k+1}.
\label{eq:return}
\end{equation}

For a policy $\pi$, the state-value and action-value functions are:

\begin{align}
V^\pi(s) &= \mathbb{E}_\pi \left[ G_t \mid s_t = s \right], \label{eq:value}\\
Q^\pi(s,a) &= \mathbb{E}_\pi \left[ G_t \mid s_t = s, a_t = a \right]. \label{eq:qvalue}
\end{align}

The Bellman expectation equation expresses the recursive structure of $V^\pi$:

\begin{equation}
V^\pi(s) =
\sum_{a \in \mathcal{A}} \pi(a \mid s)
\sum_{s' \in \mathcal{S}} P(s' \mid s,a)
\left[R(s,a) + \gamma V^\pi(s')\right].
\label{eq:bellman-expectation}
\end{equation}

\section{Temporal-difference learning}

Temporal-difference learning updates value estimates by comparing current predictions with a bootstrapped target. The one-step temporal-difference error is:

\begin{equation}
\delta_t = r_{t+1} + \gamma V(s_{t+1}) - V(s_t).
\label{eq:td-error}
\end{equation}

A tabular value estimate can then be updated as:

\begin{equation}
V(s_t) \leftarrow V(s_t) + \alpha \delta_t,
\label{eq:td-update}
\end{equation}

where $\alpha$ is the learning rate. For action-value learning, the Q-learning update is:

\begin{equation}
Q(s_t,a_t) \leftarrow Q(s_t,a_t) +
\alpha \left[
r_{t+1} + \gamma \max_{a' \in \mathcal{A}} Q(s_{t+1},a') - Q(s_t,a_t)
\right].
\label{eq:q-learning}
\end{equation}

\section{Policy optimisation}

Policy-gradient methods optimise a parameterised policy by maximising the expected return:

\begin{equation}
J(\theta) = \mathbb{E}_{\tau \sim \pi_\theta}
\left[
\sum_{t=0}^{T} \gamma^t r_{t+1}
\right].
\label{eq:policy-objective}
\end{equation}

Using the likelihood-ratio identity, the gradient can be estimated as:

\begin{equation}
\nabla_\theta J(\theta) =
\mathbb{E}_{\tau \sim \pi_\theta}
\left[
\sum_{t=0}^{T}
\nabla_\theta \log \pi_\theta(a_t \mid s_t) A^\pi(s_t,a_t)
\right],
\label{eq:policy-gradient}
\end{equation}

where $A^\pi(s,a)$ is the advantage function:

\begin{equation}
A^\pi(s,a) = Q^\pi(s,a) - V^\pi(s).
\label{eq:advantage}
\end{equation}

For clipped policy optimisation, the probability ratio between the updated and previous policies is:

\begin{equation}
\rho_t(\theta) =
\frac{\pi_\theta(a_t \mid s_t)}
{\pi_{\theta_{\mathrm{old}}}(a_t \mid s_t)}.
\label{eq:ratio}
\end{equation}

The clipped surrogate objective is:

\begin{equation}
L^{\mathrm{clip}}(\theta) =
\mathbb{E}_t
\left[
\min\left(
\rho_t(\theta) A_t,\,
\operatorname{clip}\left(\rho_t(\theta), 1-\epsilon, 1+\epsilon\right) A_t
\right)
\right].
\label{eq:ppo}
\end{equation}

The final training loss can combine policy, value, and entropy terms:

\begin{equation}
\mathcal{L}(\theta) =
-L^{\mathrm{clip}}(\theta)
+ c_v \left(V_\theta(s_t) - G_t\right)^2
- c_e \mathcal{H}\left(\pi_\theta(\cdot \mid s_t)\right),
\label{eq:combined-loss}
\end{equation}

where $c_v$ and $c_e$ control the value-loss and entropy-regularisation contributions.
